\section{Unsaturated interval bundles are relative one-handle additions}
\label{sec:unsaturated}

The incoming marked-parallelity classifier rejects every annular interface,
although many such interfaces are topologically harmless.  The distinction
that matters is whether the \emph{entire boundary of a base surface} is
attached.  The following result eliminates all nonsaturated bases at once.
It is a statement about genuine embedded interval bundles; a table of
Euler characteristics does not supply its geometric hypotheses.

\subsection{The relative surface-spine lemma}

Let $B$ be a connected compact surface with nonempty boundary.  Let
$A\subseteq\partial B$ be a nonempty compact one-dimensional submanifold,
with $a$ interval components and $c$ circle components.  Endpoints are
included.  Assume that $A\ne\partial B$.  Equivalently, if $b$ denotes the
number of boundary components, then $c<b$; maximal interval runs have
already been merged at their actual source corners.  Write
\[
 m=a+c,\qquad h=a-\chi(B),\qquad
 \lambda=h-m+1=1-\chi(B)-c.
\]

\begin{lemma}[Relative surface bands]\label{lem:relative-bands}
The pair $(B,A)$ admits a handle decomposition starting with a collar of
$A$ and adding exactly $h$ surface one-handles, with no additional zero-
or two-handles.  The collar has $m$ components.  The one-handles connect
these components, so $h\geq m-1$ and $\lambda\geq0$.
\end{lemma}

\begin{proof}
Choose a triangulation in which $A$ is a boundary subcomplex.  Since $A$ is
proper, choose an edge $e$ in the interior of a complementary boundary
interval, subdividing first if necessary.  The dual graph of the triangles,
using interior triangle adjacencies, is connected.  Root a spanning tree
at the triangle incident to $e$.

Collapse that triangle across $e$.  Then process the remaining triangles
in an order in which a parent precedes its children in the dual tree,
collapsing each triangle across its edge to its parent.  Such an edge is
free at that point: the parent triangle has been removed and there is at
most one other incident triangle.  It was not the edge across which the
parent collapsed.  All chosen edges except $e$ are interior edges of the
original triangulation.  Thus no edge or vertex of $A$ is removed.
After every triangle has been processed the surviving complex is a
connected finite graph $L$ containing $A$.

In a surface, this collapse gives a relative regular-neighborhood
presentation: reverse a triangle collapse by enlarging a collar of its
surviving two-edge path.  This local enlargement is an isotopy of the
boundary away from $A$, so a regular neighborhood of $L$ relative to $A$
is homeomorphic to $B$ with its specified attachment subset.

Take small vertex discs and edge bands for $L$, while beginning with the
collar of $A$.  Any graph vertices outside that collar can be joined to
it by a forest containing no cycle after each component of $A$ has been
contracted.  More explicitly, root one tree at each component of $A$ and
extend the trees until every vertex has been reached, assigning each
previously unreached vertex to its first reaching tree.  Absorb each
extra vertex disc along its parent band, in order away from the roots.
This cancels every additional zero-handle against one one-handle without
joining two components of the initial collar.  All remaining bands are
surface one-handles; bands may be twisted when $B$ is nonorientable.
There are no two-handles.

A collar of an interval has Euler characteristic one and a collar of a
circle has Euler characteristic zero.  Every band decreases Euler
characteristic by one.  Consequently its number is
$h=\chi(A)-\chi(B)=a-\chi(B)$.  A connected result obtained from $m$
initial components requires at least $m-1$ connecting bands.
\end{proof}

The triangulation used in this proof can be very large.  The proof is an
existence argument for a terminal topological classification, and does
not promise a polynomial-size explicit band system when the genus is
binary encoded.  The algorithm below computes its invariant count
without constructing those bands.

\subsection{A geometric relative-handle theorem}

\begin{theorem}[Proper-boundary bundle attachment]\label{thm:proper-bundle}
Let $E\to B$ be an interval bundle with orientable total space, and suppose
$B,A$ satisfy Lemma~\ref{lem:relative-bands}.  Let $K$ be a compact
three-manifold, and glue $E|_A$ into $\partial K$ by an embedding with the
usual compatible corner collars.  Suppose the resulting three-manifold
$X=K\cup E$ is orientable.

After absorbing a collar of the interface into $K$, the pair has a relative
handle decomposition consisting of exactly $a-\chi(B)$ three-dimensional
one-handles and no handles of other indices.  In particular, the assertion
holds for arbitrary embeddings of the annular components of $E|_A$ into
$\partial K$; no primitivity, slope, or null-homotopy assumption is needed.
\end{theorem}

\begin{proof}
Over a collar of $A$, the total bundle is a collar of $E|_A$.  This can be
seen separately on interval and circle components.  The interval case is
a product.  A boundary circle has an annular collar in $B$, and the
orientation character of $B$ is zero around it.  Orientability of $E$
therefore forces trivial fibre reversal there.  Thus the total bundle
over this collar is $(S^1\times I)\times I$, with the prescribed attaching
annulus as one collar face.  Gluing all these collar pieces to $K$
only enlarges its boundary collar, giving a manifold homeomorphic to $K$.
This conclusion does not depend on how the interface annulus sits on that
boundary.

Now lift the relative surface decomposition from
Lemma~\ref{lem:relative-bands}.  An individual band is a disc, so its
interval bundle is trivial.  Its lift is $D^1\times D^2$, attached along
the two discs $\partial D^1\times D^2$.  It is therefore a genuine
three-dimensional one-handle.  The attaching discs have disjoint collars
at each stage, as supplied by the actual base handle decomposition and
bundle charts.  There are exactly $a-\chi(B)$ such bands.

When a surface band is twisted, the fibre transition can reverse also;
the total space remains orientable by hypothesis.  These transitions
alter attaching parametrizations, not the fact that the lifted piece is
a three-dimensional one-handle.  The argument is consequently valid for
nonorientable bases as well.
\end{proof}

One must distinguish this theorem from a claim that all annular sums of
handlebodies are handlebodies.  Properness of $A$ is essential.  For an
unsaturated annulus base, one attached boundary circle gives $h=0$, so the
solid-torus bundle is just an outward collar of the annular interface.
For a saturated M\"obius base, the vertical annulus represents the square
of a solid-torus generator and the theorem does not apply.

\subsection{Exact free products and handlebody classification}

Suppose that $K=\bigsqcup_jK_j$, and that pairwise disjoint bundles
$E_i\to B_i$ are attached along $E_i|_{A_i}$.  Assume the complete
intersection contract: all intersections with the core are precisely
these vertical interfaces, with no unrecorded triple intersections.
Attached bases have $A_i\ne\varnothing$ and actual multiplicity one.
Assume they are all unsaturated.  Let $\Gamma$ be the bipartite multigraph
with a core vertex for every $K_j$, a base vertex for every attached $B_i$,
and one edge for each connected component of $A_i$.  Parallel edges are
retained.  Put
\[
 \lambda_i=1-\chi(B_i)-c_i.
\]
A connected component $C$ of $\Gamma$ has cycle rank
$\beta_1(C)=|E(C)|-|V(C)|+1$.

\begin{theorem}[Nonabelian unsaturated condensation]\label{thm:free-condense}
For the component $X_C$ indexed by $C$, there is a free-product decomposition
\[
 \pi_1(X_C)\cong
 \mathop{*}_{K_j\in C}\pi_1(K_j)
 *F_{r_C},\qquad
 r_C=\sum_{B_i\in C}\lambda_i+\beta_1(C).
\]
Here $F_r$ denotes the free group of rank $r$.  After choices of basepoint
paths, each natural core inclusion corresponds to the inclusion of its
free factor.  In particular all core fundamental groups inject.

Moreover, $X_C$ is a handlebody if and only if every $K_j$ in $C$ is a
handlebody.  If their genera are $g_j$, its genus is
\begin{equation}\label{eq:unsaturated-genus}
 G_C=\sum_{K_j\in C}g_j+
     \sum_{B_i\in C}\bigl(1-\chi(B_i)-c_i\bigr)+\beta_1(C).
\end{equation}
\end{theorem}

\begin{proof}
Apply Theorem~\ref{thm:proper-bundle} simultaneously to all attached bases.
The result is obtained from its core pieces by
$H=\sum_i(a_i-\chi(B_i))$ one-handles.  These handles connect precisely
the core pieces indexed by $C$, because each original $B_i$ is connected
and its interfaces are nonempty.  If $q$ denotes the number of core
pieces in $C$, choose $q-1$ of the connecting one-handles forming a tree
on them.  These tree handles give the free product of the core groups.
Each of the other $H-q+1$ handles adjoins a free generator, by van Kampen.
But
\[
 H-q+1=
 \sum_i\bigl(1-\chi(B_i)-c_i\bigr)
 +\Bigl(\sum_i(a_i+c_i)-q-\#\{B_i\in C\}+1\Bigr),
\]
and the second parenthesis is $\beta_1(C)$.  This proves the group
statement with the indicated natural inclusions.

One-handles added to orientable handlebodies yield orientable handlebodies,
with genus equal to the sum of the old genera plus $H-q+1$.  This gives
\eqref{eq:unsaturated-genus}.  Conversely, cutting the resulting
handlebody along all the one-handle co-core discs recovers the core
pieces with their added collars.  Cutting a handlebody along disjoint
properly embedded discs produces handlebodies; equivalently such a disc
system extends, after disc slides, to a complete meridian system.
Hence every recovered core piece is a handlebody.
\end{proof}

This is stronger than an integral-homology decomposition.  It explains
geometrically why unsaturated bases in the incoming interface-homology
formula impose no new relations on the core.  No assertion about an
explicit transported generating system or a compressed representation of
all boundary inclusions follows from the existence of the splitting.

\begin{corollary}[Exact ball criterion in this class]\label{cor:unsat-ball}
Assume the core genera have been certified.  Then $X_C$ is a ball exactly
when all core pieces in $C$ are balls, $C$ is a tree, and every attached
base in $C$ is orientable and planar with exactly one boundary component
not attached as a complete circle.
\end{corollary}
\begin{proof}
All terms in \eqref{eq:unsaturated-genus} are nonnegative.  For an orientable
base of genus $g$ and boundary count $b$,
$\lambda=2g+b-1-c$, which vanishes under $c<b$ exactly when $g=0$ and
$b=c+1$.  For a nonorientable base with $k\geq1$ crosscaps,
$\lambda=k+b-1-c\geq1$.  The remaining summands vanish precisely as stated.
A connected genus-zero handlebody is a ball.
\end{proof}

Unattached bordered bundles are separate handlebodies of genus
$1-\chi(B)$ and can be counted with binary multiplicity.  Their components
must be included before asserting that the \emph{whole cut} is one ball.
Closed bases are outside this operator.  A single ball component among
several components does not certify the unknot.

\subsection{Localizing the remaining two-handles}

\begin{corollary}[One two-handle per saturated base]\label{cor:saturation-localize}
Keep the genuine bundle and interface hypotheses, but allow $A=\partial B$.
Set $\varepsilon=1$ in that case and $\varepsilon=0$ otherwise.  The bundle
attachment admits a relative handle decomposition with
\[
 h_1=a-\chi(B)+\varepsilon,\qquad h_2=\varepsilon,
\]
and no additional zero- or three-handles.  Thus, after all unsaturated
attachments have been reduced, at most one additional two-handle is
needed for each saturated base.
\end{corollary}
\begin{proof}
The unsaturated case is Theorem~\ref{thm:proper-bundle}.  In the saturated
case, remove the interior of a small disc $D$ from $B$.  The punctured
surface $B^\circ$ has one additional, unattached boundary circle and
$\chi(B^\circ)=\chi(B)-1$.  Its attachment set is still $A$, with $a=0$,
so its interval bundle contributes $1-\chi(B)$ one-handles.  Restoring
$D\times I$ attaches along $\partial D\times I$, exactly the attaching
annulus of one three-dimensional two-handle.
\end{proof}

For a connected attachment-graph component $C$, let $\sigma$ be its
number of saturated bases and put
\[
 r_C=\sum_{B_i\in C}(1-\chi(B_i)-c_i+\varepsilon_i)+\beta_1(C).
\]
Taking all relative one-handles first and then the $\sigma$ two-handles
also gives the exact, generally noncanonical presentation
\begin{equation}\label{eq:bounded-relators}
 \pi_1(X_C)\cong
 \left(\mathop{*}_{K_j\in C}\pi_1(K_j)*F_{r_C}\right)
 \big/\langle\!\langle R_1,\ldots,R_\sigma\rangle\!\rangle.
\end{equation}
Indeed, the one-handles give the displayed free product as in
Theorem~\ref{thm:free-condense}, and van Kampen adds one normal relation
for each two-handle.  The original core maps may cease to be injective
after imposing these relations.  The formula bounds the number of new
relators independently of numerical base genus; it does not furnish
their words or their embedding data.

This localizes the topological source of new relations, but does not
bound the length of the surviving attaching curve.  A high-genus saturated
base can contribute a long product of commutators or squares.  An
algorithm using this corollary must retain that curve and its embedding;
its binary genus counter is not an attaching-map certificate.  The
native prototype deliberately leaves these saturated cases unresolved.

\subsection{A saturated counterexample with identical first Betti number}

Let $K$ be a solid torus with generator $x$.  In its boundary choose an
embedded annulus whose core is a primitive torus slope with longitudinal
winding number $3$.  Let $E$ be the orientable interval bundle over a
M\"obius band.  Then $E$ is a solid torus with generator $y$, while the
core of its vertical annulus represents $y^2$.  Glue these two annuli by
an orientation-compatible homeomorphism.  Van Kampen gives
\[
 \pi_1(K\cup E)=\langle x,y\mid x^3=y^2\rangle.
\]
Its abelianization is $\mathbb Z$, since $\gcd(3,2)=1$.  Its group is
nonabelian: the assignments $x\mapsto(123)$ and $y\mapsto(12)$ give a
surjection onto $S_3$.  The resulting connected manifold is consequently
not a genus-one handlebody despite having its first integral homology.
The M\"obius base is saturated, as required for exclusion from the theorem.

For comparison, use an annulus base and attach only one boundary circle
to the same winding-three annulus in $K$.  Its group presentation is
$\langle x,y\mid x^3=y\rangle\cong\mathbb Z$, and the relative collar
proof identifies the entire resulting manifold with $K$.  The positive
conclusion here comes from the collar geometry, not just the group.

\subsection{Encoded algorithm, certification, and limits}

The supplied finite operator records each core genus, each base signature
$(\chi,b,\text{orientation},\text{multiplicity})$, and the actual component
owner of every maximal interval or circle interface.  It rejects closed
bases and returns \texttt{INCONCLUSIVE} whenever an attached base is
saturated.  It then forms $\Gamma$, computes its connected components,
and applies \eqref{eq:unsaturated-genus}; isolated profiles are added as
binary counts.  A separate checker uses a core-only graph traversal and
the independent relative-handle formula
\[
 G_C=\sum_jg_j+\sum_i(a_i-\chi(B_i))-q+1.
\]
It verifies the source fingerprint and exact integer types.  Producer
functions are not invoked during replay.

For $s$ explicitly named pieces, $p$ interfaces, and maximal numerical
bit length $L$, standard graph traversal and binary arithmetic give
polynomial bit cost in $s+p+L$.  A straightforward implementation can
rescan the interface list for each component, giving an adequate
$O((s+p)^2)$ structural bound; indexing interfaces by component makes that
phase linear.  Multiplication of an isolated multiplicity by a genus is
charged at its binary cost.  Neither numerical genus nor multiplicity
is ever an iteration bound.

The native adapter invokes the incoming source-parallelity, ordered
boundary-word, and core-handlebody certificate producers.  Its verifier
replays those source proofs before deriving the finite model.  It checks
that ports are partitioned exactly once and that all ports in one
connected interface have one actual core owner and one actual base
owner.  Thus a user-supplied incidence table alone does not become a
source-normal-cut certificate.

The output is a terminal \emph{unmarked} topological census.  The proof
provides a homeomorphism and a relative band system by existence, but the
implementation does not emit their source charts.  Subsequent hierarchy
operations needing transported curves, intersection sequences, or a
lifted cutting disc require extra representations and bounds.  Saturated
planar and nonorientable bases still require genuine two-handle or other
attachment information.  This result removes a concrete transition
obstruction; it does not establish a complete quasi-polynomial recognizer.
